\section{La problématique et les objectifs}

\begin{frame}{Problématique~: Les limites des LLM en clinique}
  L'intégration de l'IA en pratique clinique offre un potentiel pour l'analyse
  des dossiers médicaux~\cite{ALGHAMDI2025102528,
  zhou2024surveylargelanguagemodels}.

  \medskip{}

  \begin{itemize}
    \item<2-> \textbf{Opacité~:} Difficulté à tracer l'origine.
    \item<3-> \textbf{Hallucinations~:} Risque de génération de données
      médicales erronées ou non vérifiées~\cite{10.1145/3703155}.
    \item<4-> \textbf{Limites du RAG classique~:} Le RAG
      classique~\cite{gao2024modularragtransformingrag} repose sur du texte
      non structuré, ce qui peut manquer de précision relationnelle.
  \end{itemize}
\end{frame}

\begin{frame}{Objectifs de \emph{Med-KAG}}
  Med-KAG\footnote{\url{https://github.com/c2fc2f/Med-KAG}} propose de passer
  d'un RAG classique à une architecture de \textbf{Knowledge Augmented
  Generation (KAG)}.

  \medskip{}

  \pause[]

  \textbf{Les deux objectifs principaux~:}
  \begin{enumerate}
    \item<2-> \textbf{Minimiser les hallucinations~:} Contraindre les réponses
      du modèle par des données factuelles issues d'une base structurée.
    \item<3-> \textbf{Fournir une justification~:} Rendre le système
      transparent en traçant l'origine du raisonnement via des connexions
      vérifiables.
  \end{enumerate}
\end{frame}

